% Options for packages loaded elsewhere
\PassOptionsToPackage{unicode}{hyperref}
\PassOptionsToPackage{hyphens}{url}
\PassOptionsToPackage{dvipsnames,svgnames,x11names}{xcolor}
\documentclass[
  a4paper,
]{article}
\usepackage{xcolor}
\usepackage[margin=20mm]{geometry}
\usepackage{amsmath,amssymb}
\setcounter{secnumdepth}{-\maxdimen} % remove section numbering
\usepackage{iftex}
\ifPDFTeX
  \usepackage[T1]{fontenc}
  \usepackage[utf8]{inputenc}
  \usepackage{textcomp} % provide euro and other symbols
\else % if luatex or xetex
  \usepackage{unicode-math} % this also loads fontspec
  \defaultfontfeatures{Scale=MatchLowercase}
  \defaultfontfeatures[\rmfamily]{Ligatures=TeX,Scale=1}
\fi
\usepackage{lmodern}
\ifPDFTeX\else
  % xetex/luatex font selection
  \setmainfont[]{DejaVu Serif}
  \setsansfont[]{DejaVu Sans}
  \setmonofont[]{DejaVu Sans Mono}
\fi
% Use upquote if available, for straight quotes in verbatim environments
\IfFileExists{upquote.sty}{\usepackage{upquote}}{}
\IfFileExists{microtype.sty}{% use microtype if available
  \usepackage[]{microtype}
  \UseMicrotypeSet[protrusion]{basicmath} % disable protrusion for tt fonts
}{}
\makeatletter
\@ifundefined{KOMAClassName}{% if non-KOMA class
  \IfFileExists{parskip.sty}{%
    \usepackage{parskip}
  }{% else
    \setlength{\parindent}{0pt}
    \setlength{\parskip}{6pt plus 2pt minus 1pt}}
}{% if KOMA class
  \KOMAoptions{parskip=half}}
\makeatother
\usepackage{longtable,booktabs,array}
\newcounter{none} % for unnumbered tables
\usepackage{calc} % for calculating minipage widths
% Correct order of tables after \paragraph or \subparagraph
\usepackage{etoolbox}
\makeatletter
\patchcmd\longtable{\par}{\if@noskipsec\mbox{}\fi\par}{}{}
\makeatother
% Allow footnotes in longtable head/foot
\IfFileExists{footnotehyper.sty}{\usepackage{footnotehyper}}{\usepackage{footnote}}
\makesavenoteenv{longtable}
\usepackage{graphicx}
\makeatletter
\newsavebox\pandoc@box
\newcommand*\pandocbounded[1]{% scales image to fit in text height/width
  \sbox\pandoc@box{#1}%
  \Gscale@div\@tempa{\textheight}{\dimexpr\ht\pandoc@box+\dp\pandoc@box\relax}%
  \Gscale@div\@tempb{\linewidth}{\wd\pandoc@box}%
  \ifdim\@tempb\p@<\@tempa\p@\let\@tempa\@tempb\fi% select the smaller of both
  \ifdim\@tempa\p@<\p@\scalebox{\@tempa}{\usebox\pandoc@box}%
  \else\usebox{\pandoc@box}%
  \fi%
}
% Set default figure placement to htbp
\def\fps@figure{htbp}
\makeatother
\setlength{\emergencystretch}{3em} % prevent overfull lines
\providecommand{\tightlist}{%
  \setlength{\itemsep}{0pt}\setlength{\parskip}{0pt}}
\usepackage{luatexja}
\usepackage{xurl}
\usepackage{float}
\usepackage{bookmark}
\IfFileExists{xurl.sty}{\usepackage{xurl}
\usepackage{float}}{} % add URL line breaks if available
\urlstyle{same}
\hypersetup{
  pdftitle={Unified Connection Hypothesis for Gravity{,} Electromagnetism{,} QCD{,} and Dirac Matter in the 3D Triangle-Loop Product (Part II)},
  pdfauthor={Noriaki Kihara (WF System Co., Ltd.) ORCID 0009-0004-6753-4020},
  colorlinks=true,
  linkcolor={Maroon},
  filecolor={Maroon},
  citecolor={Blue},
  urlcolor={Blue},
  pdfcreator={LaTeX via pandoc}}

\title{Unified Connection Hypothesis for Gravity, Electromagnetism, QCD,
and Dirac Matter in the 3D Triangle-Loop Product (Part II)}
\usepackage{etoolbox}
\makeatletter
\providecommand{\subtitle}[1]{% add subtitle to \maketitle
  \apptocmd{\@title}{\par {\large #1 \par}}{}{}
}
\makeatother
\subtitle{― External Spacetime \(O_8^{\mathrm{ext}}\) × Internal Gauge
\(O_8^{\mathrm{int}}\) Separation, Maxwell/Yang--Mills Splitting, and
Selective One-Generation Couplings ―}
\author{Noriaki Kihara (WF System Co., Ltd.) ORCID 0009-0004-6753-4020}
\date{November 28, 2026}

\begin{document}
\maketitle

{
\setcounter{tocdepth}{2}
\tableofcontents
}
\textbf{Author:} Noriaki Kihara (WF System Co., Ltd.)\\
\textbf{ORCID:} 0009-0004-6753-4020\\
\textbf{Date:} November 28, 2026\\
\textbf{Status:} Unified Hypothesis Sequel (Part II) to the
hypothesis-free Full Paper Part I (DOI:
\texttt{10.5281/zenodo.24622201}, synthesizing Rapid Communications 1--7
{[}S1--S7{]}). Directly connects the open questions of the author's
prior octonion Yang--Mills/Dirac hypothesis paper {[}K7{]} (Concept DOI:
\texttt{10.5281/zenodo.24387753}) and octonion QCD paper {[}K6{]}
(\texttt{10.5281/zenodo.24369465}) to Part I, supported by new
deterministic machine-precision verification
(\texttt{tri\_loop\_part2\_unified\_forces\_verify\_v0.py}).

\begin{center}\rule{0.5\linewidth}{0.5pt}\end{center}

\subsection{Abstract}\label{abstract}

Part I {[}P1{]} established, without external postulates, seven
structural pillars from the purely real 3-tensor triangle-loop product
\(\mathrm{Tri}(T)_{ijk} = \sum_{p,q,r} T_{ipq} T_{qjr} T_{rpk}\):
Hurwitz crystal selection, four-stage doubling-entanglement cascades,
\(SO(3)\) spin selection and \(G_2/SU(3)\) gauge counts, four emergent
quantum theorems, 4-class spherical metric distortion and critical
collapse, Hurwitz 4-stage mass hierarchies with exponential nesting, and
28D \(Spin(8)\) triality spin-\(1/2\) with CAR fermion charges
(\texttt{1:3:3:1}).

Dynamically unifying gravity, electromagnetism (Maxwell), strong
interactions (QCD Yang--Mills), and Dirac matter without double-using
octonionic axes requires an explicit \textbf{external spacetime
vs.~internal gauge separation hypothesis} consistent with the
Coleman--Mandula theorem {[}E13{]}. Here we resolve the open structural
questions of {[}K7{]} (anonymous update map, triality, and the two-ruler
dilemma) via Part I and verify four foundational pillars at machine
precision (zero randomness):

\begin{enumerate}
\def\labelenumi{\arabic{enumi}.}
\tightlist
\item
  \textbf{External \(O_8^{\mathrm{ext}}\) × Internal
  \(O_8^{\mathrm{int}}\) Separation (Coleman--Mandula Safe)}: On
  \(N = 64\) (\(O_8^{\mathrm{ext}} \otimes O_8^{\mathrm{int}}\)), all
  \(3 \times 9 = 27\) commutators between outer \(SO(3)\) spatial
  generators \(G_{\mathrm{ext}}\otimes I_8\) and inner
  \(SU(3)\times U(1)_Q\) gauge generators
  \(I_8\otimes G_{\mathrm{int}}\) vanish strictly
  (\textbf{\texttt{0.000e+00}}). By tensor-product homomorphism
  \(\mathrm{Tri}(A\otimes B) = \mathrm{Tri}(A)\otimes\mathrm{Tri}(B)\)
  (\texttt{2.86e-16}), both sectors evolve independently while sharing a
  single global clock via Layer-3 single-complex locking.
\item
  \textbf{Abelian Maxwell vs.~Non-Abelian QCD Yang--Mills from One
  Triple}: Evaluating the triple closure deficit of {[}K7{]} yields
  exact zero (\texttt{0.000e+00}) along the Abelian \(U(1)_Q\) (\(e_7\))
  direction (Maxwell type), whereas non-Abelian \(SU(3)\) 8-gluon
  plaquettes converge strictly to the \(O(\varepsilon^4)\) Yang--Mills
  action (\texttt{0.9955\ \$\textbackslash{}to\$\ 0.99993} for
  \(\varepsilon = 0.2 \to 0.025\)), with all 56 \(SU(3)\) Jacobi/Bianchi
  identities holding to \textbf{\texttt{0.000e+00}}.
\item
  \textbf{Dirac Dispersion and Exact Standard-Model Selective
  Couplings}: The spatial Clifford Dirac operator
  \(\not{D} = \sum_{a=1}^3 L_a p_a\) satisfies
  \(\|(i\not{D}+m)(i\not{D}-m) - (|p|^2-m^2)I\| = 2.591\times 10^{-15}\),
  and the triple Dirac readout
  \(\mathrm{Re}(o_A o_B o_C) = \bar{o}_A^T L(o_B) o_C\) holds to
  \texttt{1.110e-16}. Acting with \(U(1)_Q\) photon (\(\theta=0.73\))
  and \(SU(3)\) gluon (\(\theta_g=0.51\)) rotations on the 8 CAR
  one-generation states reproduces exact Standard Model selection rules
  without labels: \(\nu\) (EM=\texttt{0.000}, QCD=\texttt{0.000}),
  \(\bar{d}\) (EM=\texttt{0.243}, QCD=\texttt{0.504}), \(u\)
  (EM=\texttt{0.482}, QCD=\texttt{0.580}), and \(e^+\)
  (EM=\texttt{0.714}, QCD=\texttt{0.000}).
\item
  \textbf{Resolution of the Two-Ruler Dilemma {[}K7{]} and Gravity
  Unification}: The two candidate gravitational rulers of {[}K7{]}
  §19.5---\((1+3)+(1+3)\) spatial splitting and a scalar clock
  wave---coexist as a two-layer architecture in \(H = O_8 \otimes M_8\)
  (outer 4-class spatial metric distortion + inner \(M_8\) complex clock
  with spin-2 dragging). Open continuum targets (\(1/r^2\) Coulomb,
  \(1/r\) gravity, light-cone \(c\), and \(\alpha\)) are explicitly
  reserved for large-lattice \(O_8^{\otimes K}\) tests.
\end{enumerate}

\begin{center}\rule{0.5\linewidth}{0.5pt}\end{center}

\subsection{1. Introduction: Bridging {[}K6, K7{]}, Part I {[}P1{]}, and
Part II}\label{introduction-bridging-k6-k7-part-i-p1-and-part-ii}

In {[}K6, K7{]}, we expressed QCD in octonions and hypothesized that a
single triple product \(\mathrm{Re}(o_A o_B o_C)\) encodes Yang--Mills,
Dirac, distance, and coupling constants simultaneously, leaving open the
anonymous update map, 8D triality, and the two-ruler dilemma ({[}K7{]}
§19.3, §19.5). Part I {[}P1, S1--S7{]} resolved these mathematically via
\(\mathrm{Tri}(T)_{ijk}\). This sequel (Part II) explicitly introduces
the spacetime--gauge separation hypothesis to unify gravity,
electromagnetism, QCD, and Dirac matter, strictly tagging each claim as
\textbf{{[}Known{]}}, \textbf{{[}Established in Part I{]}},
\textbf{{[}Verified in Part II{]}}, or \textbf{{[}Hypothesis / Open{]}}.

\begin{center}\rule{0.5\linewidth}{0.5pt}\end{center}

\subsection{2. Resolving Open Questions of {[}K7{]} via Part
I}\label{resolving-open-questions-of-k7-via-part-i}

{\def\LTcaptype{none} % do not increment counter
\begin{longtable}[]{@{}
  >{\raggedright\arraybackslash}p{(\linewidth - 4\tabcolsep) * \real{0.3333}}
  >{\raggedright\arraybackslash}p{(\linewidth - 4\tabcolsep) * \real{0.3333}}
  >{\raggedright\arraybackslash}p{(\linewidth - 4\tabcolsep) * \real{0.3333}}@{}}
\toprule\noalign{}
\begin{minipage}[b]{\linewidth}\raggedright
Open Problem in 2026-11-20 Paper {[}K7{]}
\end{minipage} & \begin{minipage}[b]{\linewidth}\raggedright
Resolution via Part I {[}P1{]} \& Part II
\end{minipage} & \begin{minipage}[b]{\linewidth}\raggedright
Status
\end{minipage} \\
\midrule\noalign{}
\endhead
\bottomrule\noalign{}
\endlastfoot
Anonymous update map \(\Psi\) without external \(W_{ABC}\) avoiding
collapse to \(P^3\) ({[}K7{]} §9) & Triangle-loop product
\(\mathrm{Tri}(T)_{ijk}\) preserves anonymity and selects \(O_8\) &
{[}Established in Part I{]} \\
8D Triality exchanging Yang--Mills and Dirac readouts ({[}K7{]} §7) &
28D \(Spin(8)\) triality and spin-\(1/2\) sign flip (\(-I\)) verified on
\(O_8\) {[}S7{]} & {[}Established in Part I{]} \\
\(SU(3)\) color, electric charges, chirality, and mass hierarchy
({[}K7{]} §19.3) & \(e_7\)-fixed CAR \texttt{1:3:3:1} charges, \(SU(3)\)
Casimirs, \(M_8\) chiral split, Hurwitz hierarchy {[}S6, S7{]} &
{[}Established in Part I{]} \\
Finite resolution and duality ({[}K7{]} §17, §18) & Exact saturation of
\(\mathbb{Z}_2^3\) Donoho--Stark / Shannon bounds (\(\ln 8\)) {[}S4{]} &
{[}Established in Part I{]} \\
Two gravitational ruler hypotheses: \((1+3)+(1+3)\) vs.~scalar wave
({[}K7{]} §19.5) & Realized simultaneously in \(H = O_8 \otimes M_8\)
(outer 4-class metric + inner \(M_8\) clock) & {[}Part I{]} +
{[}Hypothesis{]} \\
\end{longtable}
}

\begin{center}\rule{0.5\linewidth}{0.5pt}\end{center}

\subsection{\texorpdfstring{3. External \(O_8^{\mathrm{ext}}\) ×
Internal \(O_8^{\mathrm{int}}\) Separation
Hypothesis}{3. External O\_8\^{}\{\textbackslash mathrm\{ext\}\} × Internal O\_8\^{}\{\textbackslash mathrm\{int\}\} Separation Hypothesis}}\label{external-o_8mathrmext-internal-o_8mathrmint-separation-hypothesis}

\textbf{Hypothesis U1 (Minimal Coleman--Mandula Separation)}: On
\(N = 64\) (\(O_8^{\mathrm{ext}} \otimes O_8^{\mathrm{int}}\)), outer
\(O_8^{\mathrm{ext}}\) carries external 3D space (\(SO(3)\), 4-class
metric distortion, discrete curvature \(\rho\), Dirac spatial operator
\(\not{D}\)), while inner \(O_8^{\mathrm{int}}\) (and its condensate
\(M_8\)) carries internal gauge symmetries (\(SU(3)\times U(1)_Q\), CAR
fermions, complex clock \(J\)).

\textbf{{[}Verified in Part II: P2-1{]}}: All \(3 \times 9 = 27\)
commutators
\([G_{\mathrm{ext}}(SO(3))\otimes I_8,\; I_8\otimes G_{\mathrm{int}}(SU(3)\times U(1)_Q)]\)
vanish to \textbf{\texttt{0.000e+00}}. By
\(\mathrm{Tri}(A\otimes B) = \mathrm{Tri}(A)\otimes\mathrm{Tri}(B)\)
(\texttt{2.86e-16}), both sectors evolve independently without violating
Coleman--Mandula {[}E13{]}, while Layer-3 single-complex locking
{[}S2{]} couples them to a shared global clock {[}Hypothesis{]}.

\begin{center}\rule{0.5\linewidth}{0.5pt}\end{center}

\subsection{4. Maxwell (Abelian) vs.~QCD Yang--Mills (Non-Abelian) from
One
Triple}\label{maxwell-abelian-vs.-qcd-yangmills-non-abelian-from-one-triple}

\textbf{{[}Verified in Part II: P2-2{]}} (Fig. 1(a)): 1.
\textbf{Electromagnetism (\(U(1)_Q\) Maxwell)}: Along the Abelian
\(e_7\) axis with \(c = -a-b\), the non-Abelian deficit
\(1 - \mathrm{Re}(o_A o_B o_C)\) vanishes identically
(\textbf{\texttt{0.000e+00}}). 2. \textbf{Strong Force (\(SU(3)\)
Yang--Mills)}: Along non-commuting \(SU(3)\) gluon generators
(\(i\Lambda_1, i\Lambda_4\)), the triple plaquette
\(1 - \frac{1}{8}\mathrm{Re}\,\mathrm{Tr}(U_A U_B U_C)\) converges
strictly to the \(O(\varepsilon^4)\) Yang--Mills limit
(\texttt{0.995540\ \$\textbackslash{}to\$\ 0.998888\ \$\textbackslash{}to\$\ 0.999722\ \$\textbackslash{}to\$\ 0.999931}
as \(\varepsilon = 0.2 \to 0.025\)). 3. \textbf{Bianchi / Jacobi
Identity}: All 56 \(SU(3)\) triples satisfy the Jacobi/Bianchi identity
to \textbf{\texttt{0.000e+00}}.

\begin{center}\rule{0.5\linewidth}{0.5pt}\end{center}

\subsection{5. Dirac Operator and Selective One-Generation Gauge
Couplings}\label{dirac-operator-and-selective-one-generation-gauge-couplings}

\textbf{{[}Verified in Part II: P2-3, P2-4{]}}: The spatial Clifford
Dirac operator \(\not{D} = \sum_{a=1}^3 L_a p_a\) obeys
\(\|(i\not{D}+m)(i\not{D}-m) - (|p|^2-m^2)I\| = \mathbf{2.591\times 10^{-15}}\),
and the triple Dirac readout
\(\mathrm{Re}(o_A o_B o_C) = \bar{o}_A^T L(o_B) o_C\) holds to
\textbf{\texttt{1.110\textbackslash{}times\ 10\^{}\{-16\}}}. Acting on
the 8 CAR one-generation states with photon (\(\theta = 0.73\)) and
8-gluon (\(\theta_g = 0.51\)) rotations yields Table 1 and Fig. 1(b):

{\def\LTcaptype{none} % do not increment counter
\begin{longtable}[]{@{}
  >{\raggedright\arraybackslash}p{(\linewidth - 8\tabcolsep) * \real{0.1667}}
  >{\centering\arraybackslash}p{(\linewidth - 8\tabcolsep) * \real{0.2083}}
  >{\centering\arraybackslash}p{(\linewidth - 8\tabcolsep) * \real{0.2083}}
  >{\centering\arraybackslash}p{(\linewidth - 8\tabcolsep) * \real{0.2083}}
  >{\centering\arraybackslash}p{(\linewidth - 8\tabcolsep) * \real{0.2083}}@{}}
\toprule\noalign{}
\begin{minipage}[b]{\linewidth}\raggedright
One-Generation Sector (8 States)
\end{minipage} & \begin{minipage}[b]{\linewidth}\centering
Charge \(Q\)
\end{minipage} & \begin{minipage}[b]{\linewidth}\centering
Phase Error
\end{minipage} & \begin{minipage}[b]{\linewidth}\centering
\textbf{EM Coupling \(\|U_Q\psi - \psi\|\)}
\end{minipage} & \begin{minipage}[b]{\linewidth}\centering
\textbf{QCD Coupling \(\max_A\|U_g\psi - \psi\|\)}
\end{minipage} \\
\midrule\noalign{}
\endhead
\bottomrule\noalign{}
\endlastfoot
\textbf{\(\nu\) (Neutrino singlet)} & \texttt{0} & \texttt{1.570e-16} &
\textbf{\texttt{0.000000} (decoupled)} & \textbf{\texttt{0.000000}
(decoupled)} \\
\textbf{\(\bar{d}\) (Anti-down triplet)} & \texttt{+1/3} &
\texttt{1.943e-16} & \textbf{\texttt{0.242733} (coupled)} &
\textbf{\texttt{0.504491} (coupled)} \\
\textbf{\(u\) (Up triplet)} & \texttt{+2/3} & \texttt{2.289e-16} &
\textbf{\texttt{0.481878} (coupled)} & \textbf{\texttt{0.580425}
(coupled)} \\
\textbf{\(e^+\) (Positron singlet)} & \texttt{+1} & \texttt{1.570e-16} &
\textbf{\texttt{0.713899} (max EM)} & \textbf{\texttt{0.000000}
(decoupled)} \\
\end{longtable}
}

\begin{figure}[H]
\centering
\includegraphics[width=0.96\linewidth,height=\textheight,keepaspectratio,alt={Part II machine-precision verification: (a) Abelian Maxwell (zero deficit) vs.~non-Abelian SU(3) Yang--Mills (O(\textbackslash varepsilon\^{}4) convergence) from one triple. (b) Exact Standard-Model selective EM and QCD couplings across all 8 one-generation states.}]{figs_en/fig_p2_ym_maxwell_and_selective_coupling_en.png}
\caption{Part II machine-precision verification: (a) Abelian Maxwell
(zero deficit) vs.~non-Abelian \(SU(3)\) Yang--Mills
(\(O(\varepsilon^4)\) convergence) from one triple. (b) Exact
Standard-Model selective EM and QCD couplings across all 8
one-generation states.}
\end{figure}

\begin{center}\rule{0.5\linewidth}{0.5pt}\end{center}

\subsection{6. Unified Architecture with Gravity and Explicit Open
Problems}\label{unified-architecture-with-gravity-and-explicit-open-problems}

Figure 2 illustrates the unified Part II architecture combining outer
gravity (4-class spherical metric distortion, discrete curvature
\(\rho\), critical collapse at \(\epsilon_c \in (0.270, 0.275)\)) with
inner electromagnetism, QCD, Dirac fermions, and the \(M_8\) complex
clock (spin-2 dragging).

\begin{figure}[H]
\centering
\includegraphics[width=0.96\linewidth,height=\textheight,keepaspectratio,alt={Part II unified architecture: separation of external spacetime O\_8\^{}\{\textbackslash mathrm\{ext\}\} (gravity, Dirac spatial derivative) and internal gauge O\_8\^{}\{\textbackslash mathrm\{int\}\}/M\_8 (EM, QCD, clock) via homomorphism, coupled by Layer-3 single-clock locking, triality, and O\_8\^{}\{\textbackslash otimes K\} nesting.}]{figs_en/fig_p2_unified_architecture_en.png}
\caption{Part II unified architecture: separation of external spacetime
\(O_8^{\mathrm{ext}}\) (gravity, Dirac spatial derivative) and internal
gauge \(O_8^{\mathrm{int}}/M_8\) (EM, QCD, clock) via homomorphism,
coupled by Layer-3 single-clock locking, triality, and
\(O_8^{\otimes K}\) nesting.}
\end{figure}

\textbf{Explicitly Unverified Targets ({[}Hypothesis / Open{]})}: (1)
Large-lattice \(O_8^{\otimes K}\) continuum tests for \(1/r^2\) Coulomb,
\(1/r\) gravity, light-cone \(c\), Lorentz factor \(\sqrt{1-v^2/c^2}\),
and light bending (factor 2); (2) dynamical electron-photon scattering
and first-principles derivation of \(\alpha \approx 1/137.036\) and
\(\alpha_s\); (3) explicit \(SU(2)_L\) weak interactions and
3-generation mass ratios.

\begin{center}\rule{0.5\linewidth}{0.5pt}\end{center}

\subsection{7. Reproducibility and Released
Files}\label{reproducibility-and-released-files}

\begin{itemize}
\tightlist
\item
  Verification script:
  \texttt{tri\_loop\_part2\_unified\_forces\_verify\_v0.py} (0.02 s)
\item
  Numerical log:
  \texttt{results/tri\_loop\_part2\_unified\_forces\_verify\_v0.txt}
\item
  Figure generator: \texttt{plot\_part2\_figures\_v0.py}
\item
  Figures (Fig. 1):
  \texttt{results/fig\_p2\_ym\_maxwell\_and\_selective\_coupling\_\{ja,en\}.png}
\item
  Figures (Fig. 2):
  \texttt{results/fig\_p2\_unified\_architecture\_\{ja,en\}.png}
\item
  Builder: \texttt{build\_tex\_part2.py}
\end{itemize}

\subsection{References}\label{references}

\subsubsection{Self-Citations}\label{self-citations}

\begin{itemize}
\tightlist
\item
  \textbf{{[}P1{]}} N. Kihara, \emph{Spontaneous Emergence of Quantum,
  Gauge, Geometric, and Fermionic Structures from the 3D Triangle-Loop
  Product (Unified Full Paper Part I)}, Zenodo (2026-11-28), Version
  DOI: \texttt{10.5281/zenodo.24622201} (Concept:
  \texttt{10.5281/zenodo.24622200}).
\item
  \textbf{{[}S1--S7{]}} N. Kihara, \emph{Rapid Communications 1--7},
  Zenodo (2026-11-28), DOIs: \texttt{10.5281/zenodo.24616931},
  \texttt{24618400}, \texttt{24619357}, \texttt{24620262},
  \texttt{24621139}, \texttt{24621454}, \texttt{24621788}.
\item
  \textbf{{[}K6{]}} N. Kihara, \emph{Representing QCD with Octonions},
  Zenodo (2026-11-19), Concept DOI: \texttt{10.5281/zenodo.24369465}.
\item
  \textbf{{[}K7{]}} N. Kihara, \emph{Reading Out Yang--Mills and Dirac
  from the Octonion Product}, Zenodo (2026-11-20), Concept DOI:
  \texttt{10.5281/zenodo.24387753}.
\end{itemize}

\subsubsection{External Primary
Literature}\label{external-primary-literature}

\begin{itemize}
\tightlist
\item
  \textbf{{[}E1{]}} C. N. Yang and R. L. Mills, \emph{Phys. Rev.}
  \textbf{96}, 191--195 (1954).
\item
  \textbf{{[}E2{]}} P. A. M. Dirac, \emph{Proc. R. Soc. Lond. A}
  \textbf{117}, 610--624 (1928).
\item
  \textbf{{[}E3{]}} K. G. Wilson, \emph{Phys. Rev.~D} \textbf{10},
  2445--2459 (1974).
\item
  \textbf{{[}E4{]}} J. C. Baez, \emph{Bull. Amer. Math. Soc.}
  \textbf{39}, 145--205 (2002).
\item
  \textbf{{[}E5{]}} M. Günaydin and F. Gürsey, \emph{J. Math. Phys.}
  \textbf{14}, 1651--1667 (1973); C. Furey, arXiv:1611.09182 (2015);
  \emph{Eur. Phys. J. C} \textbf{78}, 375 (2018).
\item
  \textbf{{[}E13{]}} S. Coleman and J. Mandula, \emph{Phys. Rev.}
  \textbf{159}, 1251--1256 (1967).
\item
  \textbf{{[}E14{]}} T. Regge, \emph{Nuovo Cimento} \textbf{19},
  558--571 (1961).
\end{itemize}

\end{document}
